\section{Alternate Hybrids}
\label{sec:hybrids}

\paraq{What does this section propose?}
Section~\ref{sec:limitations} showed that neither full split parallelism (gsplit) nor data
parallelism (DGL) is faster across the measured space. Both strategies sit at the ends of a
spectrum: gsplit computes each node once, on its owning GPU, but pays an all-to-all
\textbf{shuffle} at every layer; data parallelism never shuffles, but every GPU recomputes the
overlapping parts of its neighbourhood. In this section, we introduce two hybrids that keep gsplit's
partition-aware shuffle where it pays off and avoid or hide it where it does not: a
\emph{data-parallel hybrid}, which removes shuffles (Section~\ref{sec:hybrid-dp}), and a
\emph{pipeline hybrid}, which overlaps them with computation (Section~\ref{sec:hybrid-pipeline}).
We then measure both against gsplit (Section~\ref{sec:hybrid-empirical}).

\subsection{Hybrid Data Parallelism}
\label{sec:hybrid-dp}

\paraq{How does the data-parallel hybrid remove shuffles?}
We divide the model's layers at a \textbf{split point} $k$. Layers $0,\ldots,k{-}1$, counted
from the seed nodes, run as \textbf{data parallelism}: each GPU samples and computes the
neighbourhood of its own seed shard on replicated features, with no cross-GPU shuffle. Layers
$k,\ldots,L{-}1$ run as gsplit-style \textbf{split} parallelism, with a shuffle before each
layer (Figure~\ref{fig:data-split}): at layer $k$ each GPU still computes its own frontier but
receives that layer's inputs from their owning GPUs, and from layer $k{+}1$ on every node is
computed by its owner.
Setting $k{=}0$ recovers full split (gsplit), and $k{=}L$ recovers pure data parallelism.
% Detail: in the code k is num_redundant_layers (config.num_redundant_layer -> model num_dp and
% sampler SetNumDP); DDP already all-reduces parameter gradients, so the DP prefix needs no new
% reduction. Correctness verified by micro-tests MT1/MT2 (plans/01_dp_plus_split.md).

\paraq{What does a larger $k$ cost?}
Each layer moved into the data-parallel prefix removes one shuffle per iteration in both the
forward and backward passes. The price is \textbf{redundant computation}: neighbourhoods of
different seed shards overlap, so nodes in the overlap are computed by more than one GPU. The
layers nearest the seeds have the smallest frontiers, which makes them the cheapest to
replicate, so increasing $k$ removes shuffles in the order where replication costs least.

\begin{figure}[t]
  \centering
  % Data-split (DP-then-split) hybrid: first k layers data-parallel (no cross-GPU
% shuffle), remaining L-k layers partition-aware split (all-to-all before each layer).
% % Data-split (DP-then-split) hybrid: first k layers data-parallel (no cross-GPU
% shuffle), remaining L-k layers partition-aware split (all-to-all before each layer).
% % Data-split (DP-then-split) hybrid: first k layers data-parallel (no cross-GPU
% shuffle), remaining L-k layers partition-aware split (all-to-all before each layer).
% \input{figures/data_split} inside a figure environment.
% Layout: one row per GPU, one column per layer; each all-to-all shuffle is drawn as a
% single vertical band spanning all GPUs (instead of GPU-to-GPU arrows) to keep it legible.
\begin{tikzpicture}[
    font=\small,
    box/.style={draw, rounded corners=2pt, minimum width=1.6cm, minimum height=0.6cm, align=center},
    dp/.style={box, fill=blue!8},
    sp/.style={box, fill=red!10},
    band/.style={draw=red!60!black, fill=red!12, rounded corners=2pt},
    flow/.style={-{Latex[length=1.6mm]}, gray!70!black},
    hdr/.style={font=\scriptsize\itshape},
  ]
  \def\dy{0.85}
  % rows = GPUs; columns = layers 0..3 (k=2: layers 0,1 DP | layers 2,3 split)
  \foreach \r in {0,1,2,3}{
    \pgfmathsetmacro{\y}{-\r*\dy}
    \node[font=\scriptsize, anchor=east] at (-1.0,\y) {GPU \r};
    \node[dp] (a\r) at (0,\y)   {DP};
    \node[dp] (b\r) at (2.0,\y) {DP};
    \node[sp] (c\r) at (4.6,\y) {split};
    \node[sp] (d\r) at (7.2,\y) {split};
    \draw[flow] (a\r) -- (b\r);
  }

  % all-to-all shuffle bands (before layer 2 and before layer 3)
  \foreach \x/\n in {3.3/s1, 5.9/s2}{
    \node[band, minimum width=0.5cm, minimum height={3*\dy cm + 0.6cm}, anchor=north]
      (\n) at (\x,0.3) {};
    \node[rotate=90, font=\scriptsize, text=red!60!black] at (\n.center) {all-to-all shuffle};
  }
  \foreach \r in {0,1,2,3}{
    \draw[flow] (b\r.east) -- (b\r.east -| s1.west);
    \draw[flow] (c\r.west -| s1.east) -- (c\r.west);
    \draw[flow] (c\r.east) -- (c\r.east -| s2.west);
    \draw[flow] (d\r.west -| s2.east) -- (d\r.west);
  }

  % layer headers, then the switch point on its own row above them
  \foreach \n/\l in {a0/0, b0/1, c0/2, d0/3}{ \node[hdr, above=3pt of \n] {Layer \l}; }
  \node[font=\scriptsize, above=16pt of s1.north] (sw) {switch point $k{=}2$};
  \draw[-{Latex[length=1.6mm]}] (sw.south) -- (s1.north);

  % region braces
  \draw[decorate, decoration={brace, amplitude=4pt, mirror}]
    ($(a3.south west)+(0,-0.25)$) -- ($(b3.south east)+(0,-0.25)$)
    node[midway, below=5pt, hdr, text=blue!60!black] {data-parallel: no shuffle};
  \draw[decorate, decoration={brace, amplitude=4pt, mirror}]
    ($(s1.west |- d3.south)+(0,-0.25)$) -- ($(d3.south east)+(0,-0.25)$)
    node[midway, below=5pt, hdr, text=red!60!black] {split: shuffle before each layer};
\end{tikzpicture}
 inside a figure environment.
% Layout: one row per GPU, one column per layer; each all-to-all shuffle is drawn as a
% single vertical band spanning all GPUs (instead of GPU-to-GPU arrows) to keep it legible.
\begin{tikzpicture}[
    font=\small,
    box/.style={draw, rounded corners=2pt, minimum width=1.6cm, minimum height=0.6cm, align=center},
    dp/.style={box, fill=blue!8},
    sp/.style={box, fill=red!10},
    band/.style={draw=red!60!black, fill=red!12, rounded corners=2pt},
    flow/.style={-{Latex[length=1.6mm]}, gray!70!black},
    hdr/.style={font=\scriptsize\itshape},
  ]
  \def\dy{0.85}
  % rows = GPUs; columns = layers 0..3 (k=2: layers 0,1 DP | layers 2,3 split)
  \foreach \r in {0,1,2,3}{
    \pgfmathsetmacro{\y}{-\r*\dy}
    \node[font=\scriptsize, anchor=east] at (-1.0,\y) {GPU \r};
    \node[dp] (a\r) at (0,\y)   {DP};
    \node[dp] (b\r) at (2.0,\y) {DP};
    \node[sp] (c\r) at (4.6,\y) {split};
    \node[sp] (d\r) at (7.2,\y) {split};
    \draw[flow] (a\r) -- (b\r);
  }

  % all-to-all shuffle bands (before layer 2 and before layer 3)
  \foreach \x/\n in {3.3/s1, 5.9/s2}{
    \node[band, minimum width=0.5cm, minimum height={3*\dy cm + 0.6cm}, anchor=north]
      (\n) at (\x,0.3) {};
    \node[rotate=90, font=\scriptsize, text=red!60!black] at (\n.center) {all-to-all shuffle};
  }
  \foreach \r in {0,1,2,3}{
    \draw[flow] (b\r.east) -- (b\r.east -| s1.west);
    \draw[flow] (c\r.west -| s1.east) -- (c\r.west);
    \draw[flow] (c\r.east) -- (c\r.east -| s2.west);
    \draw[flow] (d\r.west -| s2.east) -- (d\r.west);
  }

  % layer headers, then the switch point on its own row above them
  \foreach \n/\l in {a0/0, b0/1, c0/2, d0/3}{ \node[hdr, above=3pt of \n] {Layer \l}; }
  \node[font=\scriptsize, above=16pt of s1.north] (sw) {switch point $k{=}2$};
  \draw[-{Latex[length=1.6mm]}] (sw.south) -- (s1.north);

  % region braces
  \draw[decorate, decoration={brace, amplitude=4pt, mirror}]
    ($(a3.south west)+(0,-0.25)$) -- ($(b3.south east)+(0,-0.25)$)
    node[midway, below=5pt, hdr, text=blue!60!black] {data-parallel: no shuffle};
  \draw[decorate, decoration={brace, amplitude=4pt, mirror}]
    ($(s1.west |- d3.south)+(0,-0.25)$) -- ($(d3.south east)+(0,-0.25)$)
    node[midway, below=5pt, hdr, text=red!60!black] {split: shuffle before each layer};
\end{tikzpicture}
 inside a figure environment.
% Layout: one row per GPU, one column per layer; each all-to-all shuffle is drawn as a
% single vertical band spanning all GPUs (instead of GPU-to-GPU arrows) to keep it legible.
\begin{tikzpicture}[
    font=\small,
    box/.style={draw, rounded corners=2pt, minimum width=1.6cm, minimum height=0.6cm, align=center},
    dp/.style={box, fill=blue!8},
    sp/.style={box, fill=red!10},
    band/.style={draw=red!60!black, fill=red!12, rounded corners=2pt},
    flow/.style={-{Latex[length=1.6mm]}, gray!70!black},
    hdr/.style={font=\scriptsize\itshape},
  ]
  \def\dy{0.85}
  % rows = GPUs; columns = layers 0..3 (k=2: layers 0,1 DP | layers 2,3 split)
  \foreach \r in {0,1,2,3}{
    \pgfmathsetmacro{\y}{-\r*\dy}
    \node[font=\scriptsize, anchor=east] at (-1.0,\y) {GPU \r};
    \node[dp] (a\r) at (0,\y)   {DP};
    \node[dp] (b\r) at (2.0,\y) {DP};
    \node[sp] (c\r) at (4.6,\y) {split};
    \node[sp] (d\r) at (7.2,\y) {split};
    \draw[flow] (a\r) -- (b\r);
  }

  % all-to-all shuffle bands (before layer 2 and before layer 3)
  \foreach \x/\n in {3.3/s1, 5.9/s2}{
    \node[band, minimum width=0.5cm, minimum height={3*\dy cm + 0.6cm}, anchor=north]
      (\n) at (\x,0.3) {};
    \node[rotate=90, font=\scriptsize, text=red!60!black] at (\n.center) {all-to-all shuffle};
  }
  \foreach \r in {0,1,2,3}{
    \draw[flow] (b\r.east) -- (b\r.east -| s1.west);
    \draw[flow] (c\r.west -| s1.east) -- (c\r.west);
    \draw[flow] (c\r.east) -- (c\r.east -| s2.west);
    \draw[flow] (d\r.west -| s2.east) -- (d\r.west);
  }

  % layer headers, then the switch point on its own row above them
  \foreach \n/\l in {a0/0, b0/1, c0/2, d0/3}{ \node[hdr, above=3pt of \n] {Layer \l}; }
  \node[font=\scriptsize, above=16pt of s1.north] (sw) {switch point $k{=}2$};
  \draw[-{Latex[length=1.6mm]}] (sw.south) -- (s1.north);

  % region braces
  \draw[decorate, decoration={brace, amplitude=4pt, mirror}]
    ($(a3.south west)+(0,-0.25)$) -- ($(b3.south east)+(0,-0.25)$)
    node[midway, below=5pt, hdr, text=blue!60!black] {data-parallel: no shuffle};
  \draw[decorate, decoration={brace, amplitude=4pt, mirror}]
    ($(s1.west |- d3.south)+(0,-0.25)$) -- ($(d3.south east)+(0,-0.25)$)
    node[midway, below=5pt, hdr, text=red!60!black] {split: shuffle before each layer};
\end{tikzpicture}

  \caption{\textbf{Data-parallel hybrid} with $L{=}4$ layers and split point $k{=}2$ on
  four GPUs. Layers $0$ and $1$ run data-parallel: each GPU computes its own micro-batch with
  no cross-GPU communication (blue). Layers $2$ and $3$ run split, with an all-to-all shuffle
  before each layer (red bands); from layer $3$ on every node is computed by its owning GPU. A larger $k$
  removes more shuffles but makes each GPU recompute more of the early-layer work.}
  \label{fig:data-split}
\end{figure}

\subsection{Hybrid Pipeline Parallelism}
\label{sec:hybrid-pipeline}

\paraq{How does the pipeline hybrid hide shuffles instead of removing them?}
We divide each mini-batch into two micro-batches, $A$ and $B$, which run through the same model
on two CUDA streams, staggered by one stage. While $A$ computes layer $i$, $B$ performs the
shuffle of the same layer, and the roles then alternate (Figure~\ref{fig:pipeline-split}). The
shuffle still occurs, but it overlaps with useful work, so the cost of each layer approaches
$\max(\text{compute}, \text{comm})$ instead of their sum. The two micro-batches accumulate
gradients into a single optimizer step, so the training semantics match those of one
mini-batch.
% Implementation details (plans/02_pipeline_plus_split.md): sampling of A and B is sequential
% (no overlap benefit, and it sidesteps the sampler's single-_batch constraint); the stagger
% keeps exactly one all-to-all in flight, so a single NCCL communicator suffices; A's and B's
% backward passes accumulate into one gradient, which is all-reduced once per optimizer step.

\paraq{How does it combine with the data-parallel hybrid, and what does it cost?}
We run the pipeline over full split ($k{=}0$) only, because it does not combine with a
data-parallel prefix: the prefix layers have no shuffle, so there is nothing for the other micro-batch to overlap with.
Overlap only helps when there is enough computation to cover the shuffle. Each micro-batch holds half of the
mini-batch's seeds, so the two live micro-batches together occupy roughly the memory of one
mini-batch, and an optimizer step consumes the same number of seeds as in gsplit.
% Measured runs (2026-11-14) use micro_batch_size = batch_size // 2 and include the DGL fix
% e8c6acc8 (broadcast-offset copy on the current stream), without which GAT crashed with an
% illegal memory access whenever the pipeline ran on a non-default CUDA stream.

\begin{figure}[t]
  \centering
  % Pipeline-split hybrid: 2 micro-batches (A,B) on 2 CUDA streams, staggered by one
% stage so one stream's split SHUFFLE (comm) overlaps the other's CONV (compute).
% At any instant exactly one all-to-all is in flight (single NCCL communicator).
% % Pipeline-split hybrid: 2 micro-batches (A,B) on 2 CUDA streams, staggered by one
% stage so one stream's split SHUFFLE (comm) overlaps the other's CONV (compute).
% At any instant exactly one all-to-all is in flight (single NCCL communicator).
% % Pipeline-split hybrid: 2 micro-batches (A,B) on 2 CUDA streams, staggered by one
% stage so one stream's split SHUFFLE (comm) overlaps the other's CONV (compute).
% At any instant exactly one all-to-all is in flight (single NCCL communicator).
% \input{figures/pipeline_split} inside a figure environment.
\begin{tikzpicture}[
    font=\small,
    comp/.style={draw, fill=blue!12, minimum height=0.8cm, align=center, font=\scriptsize},
    comm/.style={draw, fill=red!14, minimum height=0.8cm, align=center, font=\scriptsize},
    idle/.style={draw, dashed, fill=black!4, minimum height=0.8cm, align=center, font=\scriptsize\itshape},
    ax/.style={-{Latex}, gray},
  ]
  \def\w{1.7}   % block width = one time unit
  \def\ya{0}    % stream A row
  \def\yb{-1.1} % stream B row

  % time axis
  \draw[ax] (-0.2,-1.95) -- (7.9,-1.95) node[right, font=\scriptsize] {time};
  \foreach \t in {0,1,2,3,4}{ \draw[gray] (\t*\w,-1.9) -- (\t*\w,-2.0); }

  % row labels
  \node[left=2pt, font=\scriptsize] at (0,\ya) {stream $s_A$};
  \node[left=2pt, font=\scriptsize] at (0,\yb) {stream $s_B$};

  % Stream A: comm L0 | comp L0 | comp L1 | comp L2
  \node[comm, anchor=west, minimum width=\w cm] at (0*\w,\ya) {shuffle $L_0$};
  \node[comp, anchor=west, minimum width=\w cm] at (1*\w,\ya) {conv $L_0$};
  \node[comp, anchor=west, minimum width=\w cm] at (2*\w,\ya) {conv $L_1$};
  \node[comp, anchor=west, minimum width=\w cm] at (3*\w,\ya) {conv $L_2$};

  % Stream B: idle | shuffle L0 | conv L0 | conv L1 ...  (staggered one stage)
  \node[idle, anchor=west, minimum width=\w cm] at (0*\w,\yb) {--};
  \node[comm, anchor=west, minimum width=\w cm] at (1*\w,\yb) {shuffle $L_0$};
  \node[comp, anchor=west, minimum width=\w cm] at (2*\w,\yb) {conv $L_0$};
  \node[comp, anchor=west, minimum width=\w cm] at (3*\w,\yb) {conv $L_1$};

  % overlap highlight: A-conv aligns with B-shuffle in the same time column
  \foreach \t in {1,2,3}{
    \draw[<->, green!50!black, thick] (\t*\w+\w/2,\ya-0.45) -- (\t*\w+\w/2,\yb+0.45);
  }
  \node[font=\scriptsize\itshape, text=green!40!black, align=center] at (6.6,-0.55)
    {compute $\parallel$ comm\\(one all-to-all\\in flight)};
\end{tikzpicture}
 inside a figure environment.
\begin{tikzpicture}[
    font=\small,
    comp/.style={draw, fill=blue!12, minimum height=0.8cm, align=center, font=\scriptsize},
    comm/.style={draw, fill=red!14, minimum height=0.8cm, align=center, font=\scriptsize},
    idle/.style={draw, dashed, fill=black!4, minimum height=0.8cm, align=center, font=\scriptsize\itshape},
    ax/.style={-{Latex}, gray},
  ]
  \def\w{1.7}   % block width = one time unit
  \def\ya{0}    % stream A row
  \def\yb{-1.1} % stream B row

  % time axis
  \draw[ax] (-0.2,-1.95) -- (7.9,-1.95) node[right, font=\scriptsize] {time};
  \foreach \t in {0,1,2,3,4}{ \draw[gray] (\t*\w,-1.9) -- (\t*\w,-2.0); }

  % row labels
  \node[left=2pt, font=\scriptsize] at (0,\ya) {stream $s_A$};
  \node[left=2pt, font=\scriptsize] at (0,\yb) {stream $s_B$};

  % Stream A: comm L0 | comp L0 | comp L1 | comp L2
  \node[comm, anchor=west, minimum width=\w cm] at (0*\w,\ya) {shuffle $L_0$};
  \node[comp, anchor=west, minimum width=\w cm] at (1*\w,\ya) {conv $L_0$};
  \node[comp, anchor=west, minimum width=\w cm] at (2*\w,\ya) {conv $L_1$};
  \node[comp, anchor=west, minimum width=\w cm] at (3*\w,\ya) {conv $L_2$};

  % Stream B: idle | shuffle L0 | conv L0 | conv L1 ...  (staggered one stage)
  \node[idle, anchor=west, minimum width=\w cm] at (0*\w,\yb) {--};
  \node[comm, anchor=west, minimum width=\w cm] at (1*\w,\yb) {shuffle $L_0$};
  \node[comp, anchor=west, minimum width=\w cm] at (2*\w,\yb) {conv $L_0$};
  \node[comp, anchor=west, minimum width=\w cm] at (3*\w,\yb) {conv $L_1$};

  % overlap highlight: A-conv aligns with B-shuffle in the same time column
  \foreach \t in {1,2,3}{
    \draw[<->, green!50!black, thick] (\t*\w+\w/2,\ya-0.45) -- (\t*\w+\w/2,\yb+0.45);
  }
  \node[font=\scriptsize\itshape, text=green!40!black, align=center] at (6.6,-0.55)
    {compute $\parallel$ comm\\(one all-to-all\\in flight)};
\end{tikzpicture}
 inside a figure environment.
\begin{tikzpicture}[
    font=\small,
    comp/.style={draw, fill=blue!12, minimum height=0.8cm, align=center, font=\scriptsize},
    comm/.style={draw, fill=red!14, minimum height=0.8cm, align=center, font=\scriptsize},
    idle/.style={draw, dashed, fill=black!4, minimum height=0.8cm, align=center, font=\scriptsize\itshape},
    ax/.style={-{Latex}, gray},
  ]
  \def\w{1.7}   % block width = one time unit
  \def\ya{0}    % stream A row
  \def\yb{-1.1} % stream B row

  % time axis
  \draw[ax] (-0.2,-1.95) -- (7.9,-1.95) node[right, font=\scriptsize] {time};
  \foreach \t in {0,1,2,3,4}{ \draw[gray] (\t*\w,-1.9) -- (\t*\w,-2.0); }

  % row labels
  \node[left=2pt, font=\scriptsize] at (0,\ya) {stream $s_A$};
  \node[left=2pt, font=\scriptsize] at (0,\yb) {stream $s_B$};

  % Stream A: comm L0 | comp L0 | comp L1 | comp L2
  \node[comm, anchor=west, minimum width=\w cm] at (0*\w,\ya) {shuffle $L_0$};
  \node[comp, anchor=west, minimum width=\w cm] at (1*\w,\ya) {conv $L_0$};
  \node[comp, anchor=west, minimum width=\w cm] at (2*\w,\ya) {conv $L_1$};
  \node[comp, anchor=west, minimum width=\w cm] at (3*\w,\ya) {conv $L_2$};

  % Stream B: idle | shuffle L0 | conv L0 | conv L1 ...  (staggered one stage)
  \node[idle, anchor=west, minimum width=\w cm] at (0*\w,\yb) {--};
  \node[comm, anchor=west, minimum width=\w cm] at (1*\w,\yb) {shuffle $L_0$};
  \node[comp, anchor=west, minimum width=\w cm] at (2*\w,\yb) {conv $L_0$};
  \node[comp, anchor=west, minimum width=\w cm] at (3*\w,\yb) {conv $L_1$};

  % overlap highlight: A-conv aligns with B-shuffle in the same time column
  \foreach \t in {1,2,3}{
    \draw[<->, green!50!black, thick] (\t*\w+\w/2,\ya-0.45) -- (\t*\w+\w/2,\yb+0.45);
  }
  \node[font=\scriptsize\itshape, text=green!40!black, align=center] at (6.6,-0.55)
    {compute $\parallel$ comm\\(one all-to-all\\in flight)};
\end{tikzpicture}

  \caption{\textbf{Pipeline hybrid} for $L{=}3$ layers with two micro-batches, $A$ and $B$,
  on two CUDA streams. Each layer is a shuffle (red) followed by a conv (blue). $B$ starts one
  step after $A$, so in every overlapped step one micro-batch computes while the other
  shuffles, and exactly one all-to-all is in flight. Each step then costs
  $\max(\text{compute}, \text{comm})$ instead of their sum.}
  \label{fig:pipeline-split}
\end{figure}

\subsection{The Cost Trade-off}
\label{sec:hybrid-tradeoff}

\paraq{What does each hybrid trade against the shuffle?}
Both hybrids target the per-layer shuffle but pay for it differently. The data-parallel hybrid
replaces shuffles with redundant computation, which grows with $k$. The pipeline hybrid keeps
every shuffle and hides it behind computation, which only pays off when each micro-batch has
enough computation to cover its partner's shuffle. Whether
either pays off depends on how expensive a shuffle is relative to the computation it would
replace or hide.

\paraq{Which factors set the cost of a shuffle?}
The cost of a shuffle is roughly the feature volume moved divided by the interconnect bandwidth.
The \textbf{hidden size} sets the volume of each shuffle, the \textbf{depth} sets the number of
shuffles, and the \textbf{interconnect} sets the bandwidth. Disabling NCCL peer-to-peer transfers
to emulate PCIe lowers the effective bandwidth and raises the cost of every shuffle, so removing
or hiding shuffles should help more on PCIe than on NVLink. Since the same factors also change
the computation that the hybrids replicate or overlap, the best scheme and split point can differ
across models, graphs, and machines.

\paraq{What configuration problem do the hybrids introduce?}
The hybrids enlarge the configuration space. Each model now requires a choice of strategy---a
data-parallel hybrid with split point $k\in\{0,\ldots,L\}$, or the pipeline---and the
best choice depends jointly on the graph, the model, its depth and hidden size, and the
interconnect.

\subsection{Empirical Hybrid Sweep}
\label{sec:hybrid-empirical}

\paraq{How do we measure the hybrids?}
We measure every hybrid configuration on the same grid as Section~\ref{sec:limitations}, with the
setup of the chapter introduction, and report per-epoch and training-only times as there. For each
configuration we run the data-parallel hybrid at every interior split point
$k\in\{1,\ldots,L{-}1\}$ and the pipeline hybrid (two micro-batches of $512$ seeds each), and we
compare the fastest data-parallel variant and the pipeline against baseline gsplit ($k{=}0$).
% Detail: pipeline variants with k>0 mix the pipeline with the data-parallel hybrid, so they
% are not measured; the pipeline column isolates the effect of overlap alone.
% Detail: k = L (pure data parallelism) is not run as a gsplit variant; DGL in
% Section~\ref{sec:limitations} is the data-parallel reference.
% Detail: "best" is purely mechanical -- the minimum measured per-epoch time per family
% (gen_hybrid_table.py docstring). A family whose every variant ran out of memory is shown as OOM.
% ---------------------------------------------------------------------------------------------
% REPRODUCTION (for review; not printed)
%   Tables are auto-generated from defense_latex/results/sweep_all.csv (split rows with
%   num_redundant_layers = k, and gsplit_pipepar rows), via defense_latex/gen_hybrid_table.py;
%   do not edit tables/hybrid_*.tex by hand.
%   Counts quoted under "Observations" were derived from the rendered tables
%   (tables/hybrid_{products,papers100M}.tex) by parsing each row's Delta, winning family,
%   OOM marker, and best-DP k:
%     (recomputed 2026-11-14 from results/hybrid_ablation.csv after the pipeline re-run)
%     40 configs; hybrid faster in 37 (exceptions: GAT NVLink L=2 on both graphs, products GAT
%     NVLink L=4); Delta range -6.0% .. +38.9%; DP family fastest in 40/40; pipeline slower than
%     baseline gsplit in 40/40 (-69.9% .. -4.6%; closest: papers100M SAGE L=3 h=4096 NVLink);
%     0 OOM / crash; PCIe Delta > NVLink Delta in 20/20 pairs; SAGE Delta > GAT Delta in 19/20;
%     papers100M > products in 16/20; L=2->4 raises Delta in 7/8 ladders; h=256->4096 raises
%     Delta in 7/8; best DP k = L-1 in 37/40.
% ---------------------------------------------------------------------------------------------

\paraq{What do the tables report?}
\textbf{Results.} Tables~\ref{tab:hybrid-products} and~\ref{tab:hybrid-papers100m} report, per
configuration, baseline gsplit, the fastest data-parallel variant with its winning split point
$k$, and the pure pipeline, together with $\Delta$, the fastest hybrid's improvement over
baseline gsplit.

\FloatBarrier
% AUTO-GENERATED by defense_latex/gen_hybrid_table.py -- do not edit by hand.
\begin{table}[t]
\centering
\caption{\textbf{papers100M --- gsplit hybrids.} Per-epoch time (s), \emph{training-only} (fwd+bwd) in parentheses. \textbf{best DP}: fastest \texttt{split} with $k$ leading data-parallel layers ($k{=}1..L{-}1$); \textbf{pipeline}: pure \texttt{gsplit\_pipepar} pipeline ($k{=}0$, no redundant layers); the winning DP $k$ is shown. $\Delta$ is the fastest hybrid's improvement over baseline gsplit (\textbf{positive $=$ hybrid faster}), per-epoch $(\text{training})$. \texttt{OOM} $=$ every variant in that column ran out of memory (kept in place). Pipeline: two micro-batches of half the batch each. $0\%$ cache, 3 epochs, random features; PCIe emulated via \texttt{NCCL\_P2P\_DISABLE=1}.}
\label{tab:hybrid-papers100m}
\setlength{\tabcolsep}{4pt}
\resizebox{\textwidth}{!}{%
\begin{tabular}{|l|c|c|c|r|r|r|r|}
\hline
\textbf{Model} & $\mathbf{L}$ & $\mathbf{h}$ & \textbf{Interconnect} & \textbf{baseline gsplit} & \textbf{best DP} & \textbf{pipeline} & $\mathbf{\Delta}$ \\
\hline
\multirow{10}{*}{\textbf{GAT}} & \multirow{2}{*}{$3$} & \multirow{2}{*}{$256$} & PCIe (emul.) & 22.05 (14.77) & 20.62 (14.01)~{\footnotesize $k2$} & 34.30 (26.53) & $+6.5\%$\,(+5\%)~{\footnotesize DP} \\
& & & NVLink & 21.21 (13.52) & 20.82 (13.40)~{\footnotesize $k1$} & 34.15 (25.82) & $+1.8\%$\,(+1\%)~{\footnotesize DP} \\
\cline{2-8}
& \multirow{2}{*}{$3$} & \multirow{2}{*}{$1024$} & PCIe (emul.) & 34.89 (28.24) & 31.04 (24.11)~{\footnotesize $k2$} & 45.89 (38.58) & $+11.0\%$\,(+15\%)~{\footnotesize DP} \\
& & & NVLink & 30.24 (23.15) & 28.79 (21.81)~{\footnotesize $k2$} & 40.88 (32.90) & $+4.8\%$\,(+6\%)~{\footnotesize DP} \\
\cline{2-8}
& \multirow{2}{*}{$3$} & \multirow{2}{*}{$4096$} & PCIe (emul.) & 166.23 (158.05) & 143.88 (136.32)~{\footnotesize $k2$} & 183.15 (174.23) & $+13.4\%$\,(+14\%)~{\footnotesize DP} \\
& & & NVLink & 147.96 (139.93) & 142.38 (134.51)~{\footnotesize $k2$} & 160.43 (150.62) & $+3.8\%$\,(+4\%)~{\footnotesize DP} \\
\cline{2-8}
& \multirow{2}{*}{$2$} & \multirow{2}{*}{$1024$} & PCIe (emul.) & 12.66 (8.83) & 12.51 (9.13)~{\footnotesize $k1$} & 21.51 (17.71) & $+1.2\%$\,(-3\%)~{\footnotesize DP} \\
& & & NVLink & 13.27 (9.08) & 13.88 (10.09)~{\footnotesize $k1$} & 21.56 (17.45) & $-4.6\%$\,(-11\%)~{\footnotesize DP} \\
\cline{2-8}
& \multirow{2}{*}{$4$} & \multirow{2}{*}{$1024$} & PCIe (emul.) & 120.49 (104.27) & 101.46 (82.67)~{\footnotesize $k3$} & 134.90 (117.23) & $+15.8\%$\,(+21\%)~{\footnotesize DP} \\
& & & NVLink & 101.44 (83.82) & 95.34 (79.33)~{\footnotesize $k2$} & 110.05 (92.67) & $+6.0\%$\,(+5\%)~{\footnotesize DP} \\
\hline
\multirow{10}{*}{\textbf{SAGE}} & \multirow{2}{*}{$3$} & \multirow{2}{*}{$256$} & PCIe (emul.) & 13.86 (6.59) & 11.57 (4.86)~{\footnotesize $k2$} & 20.91 (12.91) & $+16.5\%$\,(+26\%)~{\footnotesize DP} \\
& & & NVLink & 13.36 (5.76) & 12.33 (5.04)~{\footnotesize $k2$} & 20.80 (12.44) & $+7.7\%$\,(+12\%)~{\footnotesize DP} \\
\cline{2-8}
& \multirow{2}{*}{$3$} & \multirow{2}{*}{$1024$} & PCIe (emul.) & 21.37 (14.69) & 15.02 (8.39)~{\footnotesize $k2$} & 26.38 (19.10) & $+29.7\%$\,(+43\%)~{\footnotesize DP} \\
& & & NVLink & 16.90 (9.34) & 14.31 (6.97)~{\footnotesize $k2$} & 22.07 (13.84) & $+15.3\%$\,(+25\%)~{\footnotesize DP} \\
\cline{2-8}
& \multirow{2}{*}{$3$} & \multirow{2}{*}{$4096$} & PCIe (emul.) & 78.83 (70.24) & 61.99 (53.98)~{\footnotesize $k2$} & 97.43 (88.33) & $+21.4\%$\,(+23\%)~{\footnotesize DP} \\
& & & NVLink & 62.59 (53.21) & 54.72 (46.23)~{\footnotesize $k2$} & 65.50 (55.99) & $+12.6\%$\,(+13\%)~{\footnotesize DP} \\
\cline{2-8}
& \multirow{2}{*}{$2$} & \multirow{2}{*}{$1024$} & PCIe (emul.) & 8.40 (4.60) & 6.90 (3.53)~{\footnotesize $k1$} & 12.73 (9.12) & $+17.9\%$\,(+23\%)~{\footnotesize DP} \\
& & & NVLink & 7.89 (4.29) & 7.29 (3.38)~{\footnotesize $k1$} & 13.04 (8.86) & $+7.6\%$\,(+21\%)~{\footnotesize DP} \\
\cline{2-8}
& \multirow{2}{*}{$4$} & \multirow{2}{*}{$1024$} & PCIe (emul.) & 76.93 (59.20) & 50.92 (32.77)~{\footnotesize $k3$} & 81.30 (63.64) & $+33.8\%$\,(+45\%)~{\footnotesize DP} \\
& & & NVLink & 51.12 (34.66) & 44.55 (26.64)~{\footnotesize $k3$} & 57.22 (39.66) & $+12.9\%$\,(+23\%)~{\footnotesize DP} \\
\hline
\end{tabular}
}
\par\vspace{2pt}\footnotesize\textit{Preliminary observations (to verify):}
\begin{itemize}\setlength{\itemsep}{0pt}\setlength{\parskip}{0pt}
  \item \textbf{DP hybrid (\texttt{split}, $k>0$)} beats baseline gsplit in 19 of 20 cells (exception: GAT NVLink $L{=}2$, $-4.6\%$); the largest interior $k$ ($k{=}L{-}1$) is its fastest variant in 18 of 20. Largest gains on \textbf{PCIe + deep/wide} (SAGE PCIe $L{=}4$: $+33.8\%$; $L{=}3,h{=}1024$: $+29.7\%$).
  \item \textbf{Pure pipeline} (two half-size micro-batches) is slower than baseline gsplit in all 20 cells ($-69.9\%$ to $-4.6\%$); closest at $L{=}3,h{=}4096$ on NVLink.
\end{itemize}
\end{table}

% AUTO-GENERATED by defense_latex/gen_hybrid_table.py -- do not edit by hand.
\begin{table}[t]
\centering
\caption{\textbf{products --- gsplit hybrids.} Per-epoch time (s), \emph{training-only} (fwd+bwd) in parentheses. \textbf{best DP}: fastest \texttt{split} with $k$ leading data-parallel layers ($k{=}1..L{-}1$); \textbf{pipeline}: pure \texttt{gsplit\_pipepar} pipeline ($k{=}0$, no redundant layers); the winning DP $k$ is shown. $\Delta$ is the fastest hybrid's improvement over baseline gsplit (\textbf{positive $=$ hybrid faster}), per-epoch $(\text{training})$. \texttt{OOM} $=$ every variant in that column ran out of memory (kept in place). Pipeline: two micro-batches of half the batch each. $0\%$ cache, 3 epochs, random features; PCIe emulated via \texttt{NCCL\_P2P\_DISABLE=1}.}
\label{tab:hybrid-products}
\setlength{\tabcolsep}{4pt}
\resizebox{\textwidth}{!}{%
\begin{tabular}{|l|c|c|c|r|r|r|r|}
\hline
\textbf{Model} & $\mathbf{L}$ & $\mathbf{h}$ & \textbf{Interconnect} & \textbf{baseline gsplit} & \textbf{best DP} & \textbf{pipeline} & $\mathbf{\Delta}$ \\
\hline
\multirow{10}{*}{\textbf{GAT}} & \multirow{2}{*}{$3$} & \multirow{2}{*}{$256$} & PCIe (emul.) & 7.22 (4.73) & 6.90 (4.37)~{\footnotesize $k2$} & 10.12 (6.97) & $+4.4\%$\,(+8\%)~{\footnotesize DP} \\
& & & NVLink & 6.43 (3.66) & 6.32 (3.49)~{\footnotesize $k2$} & 8.93 (5.54) & $+1.7\%$\,(+5\%)~{\footnotesize DP} \\
\cline{2-8}
& \multirow{2}{*}{$3$} & \multirow{2}{*}{$1024$} & PCIe (emul.) & 16.48 (13.96) & 14.24 (11.68)~{\footnotesize $k2$} & 19.13 (16.00) & $+13.6\%$\,(+16\%)~{\footnotesize DP} \\
& & & NVLink & 13.99 (11.18) & 13.63 (10.84)~{\footnotesize $k2$} & 16.97 (13.53) & $+2.6\%$\,(+3\%)~{\footnotesize DP} \\
\cline{2-8}
& \multirow{2}{*}{$3$} & \multirow{2}{*}{$4096$} & PCIe (emul.) & 88.47 (85.89) & 78.60 (76.00)~{\footnotesize $k2$} & 103.99 (100.69) & $+11.2\%$\,(+12\%)~{\footnotesize DP} \\
& & & NVLink & 78.69 (75.79) & 78.13 (75.18)~{\footnotesize $k2$} & 90.44 (86.89) & $+0.7\%$\,(+1\%)~{\footnotesize DP} \\
\cline{2-8}
& \multirow{2}{*}{$2$} & \multirow{2}{*}{$1024$} & PCIe (emul.) & 3.38 (2.49) & 2.89 (2.13)~{\footnotesize $k1$} & 3.99 (3.11) & $+14.5\%$\,(+14\%)~{\footnotesize DP} \\
& & & NVLink & 3.18 (2.04) & 3.37 (2.13)~{\footnotesize $k1$} & 4.23 (3.00) & $-6.0\%$\,(-4\%)~{\footnotesize DP} \\
\cline{2-8}
& \multirow{2}{*}{$4$} & \multirow{2}{*}{$1024$} & PCIe (emul.) & 70.85 (64.72) & 63.87 (56.67)~{\footnotesize $k3$} & 93.67 (84.47) & $+9.9\%$\,(+12\%)~{\footnotesize DP} \\
& & & NVLink & 55.33 (48.96) & 55.94 (49.69)~{\footnotesize $k1$} & 77.75 (68.51) & $-1.1\%$\,(-1\%)~{\footnotesize DP} \\
\hline
\multirow{10}{*}{\textbf{SAGE}} & \multirow{2}{*}{$3$} & \multirow{2}{*}{$256$} & PCIe (emul.) & 5.08 (2.56) & 4.27 (1.75)~{\footnotesize $k2$} & 6.55 (3.40) & $+15.9\%$\,(+32\%)~{\footnotesize DP} \\
& & & NVLink & 4.31 (1.40) & 4.01 (1.04)~{\footnotesize $k2$} & 5.89 (2.32) & $+7.0\%$\,(+26\%)~{\footnotesize DP} \\
\cline{2-8}
& \multirow{2}{*}{$3$} & \multirow{2}{*}{$1024$} & PCIe (emul.) & 8.24 (5.60) & 5.54 (2.96)~{\footnotesize $k2$} & 9.58 (6.45) & $+32.8\%$\,(+47\%)~{\footnotesize DP} \\
& & & NVLink & 5.84 (2.93) & 5.13 (2.27)~{\footnotesize $k2$} & 7.06 (3.58) & $+12.2\%$\,(+23\%)~{\footnotesize DP} \\
\cline{2-8}
& \multirow{2}{*}{$3$} & \multirow{2}{*}{$4096$} & PCIe (emul.) & 27.60 (24.83) & 16.87 (14.06)~{\footnotesize $k2$} & 30.45 (27.06) & $+38.9\%$\,(+43\%)~{\footnotesize DP} \\
& & & NVLink & 18.30 (15.17) & 16.65 (13.46)~{\footnotesize $k2$} & 19.89 (16.05) & $+9.0\%$\,(+11\%)~{\footnotesize DP} \\
\cline{2-8}
& \multirow{2}{*}{$2$} & \multirow{2}{*}{$1024$} & PCIe (emul.) & 1.76 (0.90) & 1.51 (0.69)~{\footnotesize $k1$} & 2.40 (1.52) & $+14.2\%$\,(+23\%)~{\footnotesize DP} \\
& & & NVLink & 1.99 (0.80) & 1.96 (0.72)~{\footnotesize $k1$} & 2.76 (1.50) & $+1.5\%$\,(+10\%)~{\footnotesize DP} \\
\cline{2-8}
& \multirow{2}{*}{$4$} & \multirow{2}{*}{$1024$} & PCIe (emul.) & 38.77 (32.61) & 25.84 (18.52)~{\footnotesize $k3$} & 48.14 (38.83) & $+33.4\%$\,(+43\%)~{\footnotesize DP} \\
& & & NVLink & 23.03 (16.59) & 21.74 (14.52)~{\footnotesize $k3$} & 29.54 (20.16) & $+5.6\%$\,(+12\%)~{\footnotesize DP} \\
\hline
\end{tabular}
}
\par\vspace{2pt}\footnotesize\textit{Preliminary observations (to verify):}
\begin{itemize}\setlength{\itemsep}{0pt}\setlength{\parskip}{0pt}
  \item \textbf{DP hybrid} beats baseline gsplit in 18 of 20 cells (exceptions: GAT NVLink $L{=}2$, $-6.0\%$; $L{=}4,h{=}1024$, $-1.1\%$); $k{=}L{-}1$ is its fastest variant in 19 of 20. Largest gains on \textbf{PCIe + deep/wide} (SAGE PCIe $L{=}3,h{=}4096$: $+38.9\%$; $L{=}4$: $+33.4\%$).
  \item \textbf{Pure pipeline} is slower than baseline gsplit in all 20 cells ($-40.5\%$ to $-8.7\%$); closest at $L{=}3,h{=}4096$ on NVLink.
\end{itemize}
\end{table}

\FloatBarrier

\paraq{What do the tables show?}
\textbf{Observations.} Across the $40$ configurations, the fastest hybrid is faster than baseline
gsplit in $37$, with a per-epoch $\Delta$ ranging from $-6\%$ to $+39\%$. The three exceptions
are GAT on NVLink: $L{=}2$ on both graphs and $L{=}4$ on products.
\begin{itemize}\itemsep1pt
  \item \emph{Strategy.} The data-parallel hybrid is the fastest hybrid in all $40$
        configurations. The pure pipeline is slower than baseline gsplit in all $40$, by $4.6\%$
        (papers100M SAGE, $L{=}3$, $h{=}4096$, NVLink) to $70\%$ (papers100M GAT, $L{=}2$,
        PCIe); within each graph and model, its gap is smallest at $L{=}3$, $h{=}4096$ on NVLink.
  \item \emph{Split point.} The fastest data-parallel variant uses the largest interior split
        point, $k{=}L{-}1$, in $37$ of the $40$ configurations.
  \item \emph{Interconnect.} The gain is larger on emulated PCIe than on NVLink in all $20$ paired
        configurations, but a hybrid still wins in $17$ of the $20$ NVLink configurations.
  \item \emph{Depth and hidden size.} Increasing $h$ from $256$ to $4096$ raises $\Delta$ in $7$
        of the $8$ hidden-size ladders; increasing $L$ from $2$ to $4$ raises it in $7$ of the $8$
        depth ladders.
  \item \emph{Model.} SAGE gains more than GAT in $19$ of the $20$ paired configurations, the
        opposite of the model trend in gsplit's gain over DGL (Section~\ref{sec:limitations}).
  \item \emph{Memory.} Every configuration of both hybrids completed; none ran out of memory.
\end{itemize}

\paraq{What does this section conclude?}
\textbf{Conclusion.} The hybrids improve on gsplit in $37$ of $40$ configurations, but the best
split point and the size of the gain vary across configurations, and the pure pipeline does not
improve on gsplit in any of them. Choosing a configuration therefore requires searching over the
strategy, the split point $k$, and the hardware. In this section, we performed that search by brute
force, measuring every variant on real multi-GPU hardware, which is too expensive to repeat
for every new model, graph, or machine. Section~\ref{sec:roofline} replaces this sweep with a
roofline that explains and ranks the strategies from one sampling pass.
